\newpage
\section{Projektowanie}

\subsection{Architektura rozwiązania}

Projekt został podzielony na kilka logicznych modułów, z których każdy odpowiada za inny fragment eksperymentu:

\begin{itemize}
    \item \texttt{dataset.py} --- przygotowanie danych, transformacje i augmentacje,
    \item \texttt{split\_dataset.py} --- przygotowanie struktury \texttt{train/val/test},
    \item \texttt{model.py} --- tworzenie modelu na bazie biblioteki \texttt{timm},
    \item \texttt{train.py} --- trening, checkpointy, metryki i early stopping,
    \item \texttt{predict.py} oraz \texttt{inference.py} --- inferencja dla pojedynczego obrazu,
    \item \texttt{robustness\_eval.py} --- test odporności na spadek jakości,
    \item \texttt{deepfake\_faces.py} oraz \texttt{face\_utils.py} --- wykrywanie twarzy i przygotowanie cropów,
    \item \texttt{compare\_face\_models.py} --- porównanie modelu globalnego i twarzowego,
    \item \texttt{error\_analysis.py} --- analiza błędów i eksport przykładów,
    \item \texttt{web\_demo.py} --- lokalny interfejs demonstracyjny.
\end{itemize}

Takie rozdzielenie upraszcza dalszy rozwój projektu, ponieważ etap globalny i etap twarzowy mogą być rozwijane niezależnie.

\subsection{Dane}

W eksperymentach wykorzystano zbiór danych zorganizowany w klasy \texttt{real} i \texttt{fake}, a następnie podzielony na części:

\begin{itemize}
    \item zbiór treningowy: 40800 obrazów,
    \item zbiór walidacyjny: 7200 obrazów,
    \item zbiór testowy: 12000 obrazów.
\end{itemize}

Taki podział pozwala oddzielić proces uczenia od strojenia i od finalnej oceny modelu.

W dalszym etapie przygotowano również:

\begin{itemize}
    \item curated zbiór portretów do treningu modelu twarzowego,
    \item mały benchmark OOD z nowymi generatorami GPT i Gemini / Nano2,
    \item małe zbiory adaptacyjne do prób dostrajania modeli.
\end{itemize}

\subsection{Model bazowy}

Jako model bazowy wybrano architekturę \texttt{ResNet18} \cite{he2016resnet} dostępną w bibliotece \texttt{timm}. Wybór ten wynikał z kilku powodów:

\begin{itemize}
    \item model jest lekki obliczeniowo i nadaje się do szybkich eksperymentów,
    \item posiada sprawdzoną architekturę dla klasyfikacji obrazów,
    \item można go uruchamiać efektywnie na GPU z wykorzystaniem wag wstępnie wytrenowanych.
\end{itemize}

Model przyjmuje obrazy przeskalowane do rozmiaru \texttt{224 x 224} i zwraca prawdopodobieństwo dla dwóch klas: \texttt{fake} oraz \texttt{real}.

\subsection{Etap twarzowy}

Do analizy portretów zastosowano osobny pipeline. Najpierw z obrazu wejściowego wykrywana jest twarz przy użyciu klasyfikatora Haar Cascade \cite{viola2001rapid}, a następnie wycinany jest crop w stylu \texttt{face} lub \texttt{portrait}. Tak przygotowany obraz trafia do osobnego klasyfikatora twarzowego opartego o \texttt{ConvNeXt Tiny} \cite{liu2022convnext}.

Takie rozdzielenie miało umożliwić:

\begin{itemize}
    \item odciążenie modelu od pełnego kontekstu kadru,
    \item skupienie analizy na obszarze najbardziej istotnym dla deepfake'ów,
    \item porównanie zachowania modelu globalnego i wyspecjalizowanego modelu twarzowego.
\end{itemize}

\subsection{Augmentacje i odporność}

Ze względu na charakter problemu już na etapie projektowania uznano, że zwykły trening na czystych obrazach nie będzie wystarczający. Dlatego zastosowano augmentacje związane z pogorszeniem jakości:

\begin{itemize}
    \item kompresję JPEG,
    \item rozmycie Gaussa,
    \item dodawanie szumu,
    \item operacje downscale-upscale.
\end{itemize}

Celem tych modyfikacji było uodpornienie modelu na część przypadków spotykanych w praktyce, zwłaszcza po przejściu obrazów przez komunikatory i serwisy internetowe.

\subsection{Środowisko uruchomieniowe}

Projekt jest rozwijany w Pythonie z wykorzystaniem między innymi bibliotek:

\begin{itemize}
    \item \texttt{PyTorch},
    \item \texttt{torchvision},
    \item \texttt{timm},
    \item \texttt{Pillow},
    \item \texttt{OpenCV},
    \item \texttt{Gradio}.
\end{itemize}

Trening został zoptymalizowany pod GPU poprzez:

\begin{itemize}
    \item wykorzystanie \texttt{AMP},
    \item automatyczny dobór \texttt{batch size},
    \item \texttt{channels\_last},
    \item \texttt{cudnn\_benchmark},
    \item dostrojenie \texttt{DataLoadera}.
\end{itemize}

Te decyzje projektowe pozwoliły wyraźnie poprawić praktyczne wykorzystanie karty graficznej podczas treningu.
